% !TeX root = ../../Book1.tex
% 纲要标题：### 17.2 乘法群与 Frobenius
% 纲要位置：计划纲要.md，第 593 行。

\section{乘法群与 Frobenius}
\label{b1:sec:1702}

% 纲要内容（仅作写作计划，尚非教材正文）：
%
% * 有限域乘法群的循环性
% * Frobenius 自同构
% * 有限域的子域与扩张；区分域生成元与乘法本原元，解释交域与合成域的次数
% * 乘法群循环性的主要引理已在第16.4节建立；本节用它计算有限域的生成元并组织 Frobenius 及子域的应用。
%

有限域的加法来自特征为 \(p\) 的向量空间，乘法却有更强的结构：全部非零元素可以由一个元素的幂列出。与此相配合，Frobenius 的迭代既描述共轭，也能直接辨认子域。

\subsection{有限域乘法群的循环性}
\label{b1:subsec:1702:01}

\begin{corollary}\label{b1:cor:finite-field-cyclic}
设 \(q\) 为素数幂。则有限域 \(\mathbb F_q\) 的乘法群是 \(q-1\) 阶循环群。
\end{corollary}
\begin{proof}
\cref{b1:lem:finite-multiplicative-cyclic}已经证明：任意域 \(K\) 的乘法群 \(K^\times\) 的每个有限子群都是循环群。取 \(K=\mathbb F_q\)、\(G=\mathbb F_q^\times\)，则 \(G\) 有 \(q-1\) 个元素，故 \(G\cong C_{q-1}\)。
\end{proof}
设 \(q=p^r\) 为素数幂，\(g\in\mathbb F_q^\times\)。如果 \(g\) 生成乘法群 \(\mathbb F_q^\times\)，则称 \(g\) 为有限域 \(\mathbb F_q\) 的\term[benyuan-yuan-youxianyu]{本原元}[primitive element of a finite field]。这里的意思强于“生成域扩张的元素”：乘法本原元 \(g\) 的幂给出全部非零元素，所以当然有 \(\mathbb F_p(g)=\mathbb F_q\)。域生成则只要求 \(g\) 不属于 \(\mathbb F_q\) 的任何真子域：若 \(\mathbb F_p(g)\) 是真子域，它本身就包含 \(g\)；反过来，含有 \(g\) 的任意子域都包含 \(\mathbb F_p(g)\)。这个条件并不要求 \(g\) 的乘法阶为 \(q-1\)。例如在 \(\mathbb F_9=\mathbb F_3(u)\)、\(u^2=-1\) 中，\(u\) 生成二次扩张，但其乘法阶为四，小于八。

\ProseParagraphBreak
循环群中，若 \(g\) 的阶为 \(q-1\)，则 \(g^a\) 生成整个群当且仅当 \(\gcd(a,q-1)=1\)。因此本原元的个数为 \(\varphi(q-1)\)，其中 \(\varphi(m)\) 表示模 \(m\) 与 \(m\) 互素的剩余类个数。

\subsection{Frobenius 自同构}
\label{b1:subsec:1702:02}

在特征 \(p\) 的域中，\(a\mapsto a^p\) 保持加法、乘法和单位，并且单射；有限域上的单射必为满射，所以它是自同构。对 \(\mathbb F_{q^n}/\mathbb F_q\)，其中 \(q=p^r\)，使用其第 \(r\) 次迭代
\[
\sigma:\mathbb F_{q^n}\longrightarrow\mathbb F_{q^n},\qquad a\longmapsto a^q.
\]
对每个 \(c\in\mathbb F_q\)，都有 \(c^q=c\)，所以 \(\sigma\) 逐点固定 \(\mathbb F_q\)，确是这个扩张的自同构，称为扩张 \(\mathbb F_{q^n}/\mathbb F_q\) 的 Frobenius 自同构。由于每个 \(a\in\mathbb F_{q^n}\) 都满足 \(a^{q^n}=a\)，有 \(\sigma^n=1\)。要排除 \(\sigma\) 的阶比 \(n\) 小，只须排除某个 \(0<d<n\) 满足 \(\sigma^d=1\)。后一个等式要求 \(a^{q^d}=a\) 对全部 \(q^n\) 个元素成立，即它们全是 \(x^{q^d}-x\) 的根；但其次数仅为 \(q^d<q^n\)，不可能有这么多根。因此 \(\sigma\) 的阶恰为 \(n\)。其不动元为 \(x^q-x\) 在 \(\mathbb F_{q^n}\) 中的全部根，恰组成 \(\mathbb F_q\)。

\begin{theorem}\label{b1:thm:finite-field-automorphisms}
设 \(q\) 为素数幂，\(n\geq1\) 为整数。则 \(\mathbb F_{q^n}\) 上保持 \(\mathbb F_q\) 中每个元素不变的自同构恰为 \(1,\sigma,\ldots,\sigma^{n-1}\)，其中 \(\sigma(a)=a^q\)。
\end{theorem}
\begin{proof}
由 \(\sigma\) 的阶为 \(n\)，所列 \(n\) 个自同构互异。为了证明不存在其他自同构，只需找一个生成域扩张的元素，再限制它的像。取乘法本原元 \(g\)，因为它的幂给出全部非零元素，便自动有 \(\mathbb F_{q^n}=\mathbb F_q(g)\)。基计数说明 \([\mathbb F_{q^n}:\mathbb F_q]=n\)，故\cref{b1:prop:minimal-polynomial-degree}给出 \(\deg m_{g,\mathbb F_q}=n\)。任一 \(\mathbb F_q\) 自同构 \(\tau\) 都由 \(\tau(g)\) 唯一决定；又因它固定最小多项式的系数，对 \(m_{g,\mathbb F_q}(g)=0\) 施以 \(\tau\)，得到 \(m_{g,\mathbb F_q}(\tau(g))=0\)。次数为 \(n\) 的多项式至多有 \(n\) 个根，所以这样的自同构至多有 \(n\) 个。这就是此前用生成元的共轭像计数嵌入的方法，已列出的 \(n\) 个自同构因而包含全部。
\end{proof}

\subsection{有限域的子域与扩张}
\label{b1:subsec:1702:03}

\begin{theorem}\label{b1:thm:finite-field-subfields}
设 \(q\) 为素数幂，\(m,n\) 为正整数，并在 \(\mathbb F_q\) 的一个固定代数闭包中考虑这些有限域。\(\mathbb F_{q^n}\) 中包含 \(\mathbb F_q\) 的子域恰为 \(\mathbb F_{q^d}\)，其中 \(d\mid n\)，每个这样的子域唯一。此外
\[
\mathbb F_{q^m}\cap\mathbb F_{q^n}=\mathbb F_{q^{\gcd(m,n)}},\qquad
\mathbb F_{q^m}\mathbb F_{q^n}=\mathbb F_{q^{\operatorname{lcm}(m,n)}}.
\]
\end{theorem}
\begin{proof}
若 \(E\) 是所述子域，记 \([E:\mathbb F_q]=d\)。\cref{b1:thm:tower-law}给出 \(d\mid n\)，而 \(|E|=q^d\)。对非零 \(a\in E\)，由 \(E^\times\) 的阶为 \(q^d-1\) 得 \(a^{q^d}=a\)，零也满足此式。因此 \(E\) 包含于固定代数闭包中的 \(x^{q^d}-x\) 的根集。根集恰有 \(q^d\) 个元素，与 \(E\) 的基数相同，故两者相等，这证明了子域的唯一性。反过来，若 \(d\mid n\)，任意满足 \(a^{q^d}=a\) 的元素在迭代 \(n/d\) 次后满足 \(a^{q^n}=a\)，所以代数闭包中的 \(\mathbb F_{q^d}\) 包含于 \(\mathbb F_{q^n}\)。

令 \(g=\gcd(m,n)\)，并记交域 \(E=\mathbb F_{q^m}\cap\mathbb F_{q^n}\) 在 \(\mathbb F_q\) 上的次数为 \(e\)。由两条域塔，\(e\mid m\) 且 \(e\mid n\)，所以 \(e\mid g\)。另一方面，\(g\mid m,n\) 使 \(\mathbb F_{q^g}\) 同时包含于两个域，因而包含于 \(E\)；塔式公式又给出 \(g\mid e\)。故 \(e=g\)，结合已证的子域唯一性，得到交域公式。

令 \(\ell=\operatorname{lcm}(m,n)\)，并记合成域 \(M=\mathbb F_{q^m}\mathbb F_{q^n}\) 在 \(\mathbb F_q\) 上的次数为 \(s\)。两个域都包含于 \(\mathbb F_{q^\ell}\)，所以由它们生成的 \(M\) 也包含于该域，塔式公式给出 \(s\mid\ell\)。另一方面，\(M\) 包含这两个域，故 \(m\mid s\) 且 \(n\mid s\)，于是 \(\ell\mid s\)。所以 \(s=\ell\)，再由子域唯一性得到合成域公式。
\end{proof}
子域定理把有限域的包含关系化为指数的整除关系：\(\mathbb F_{q^a}\subseteq\mathbb F_{q^b}\) 当且仅当 \(a\mid b\)，交域对应最大公因数，合成域对应最小公倍数。例如 \(\mathbb F_{64}\) 包含 \(\mathbb F_2,\mathbb F_4,\mathbb F_8,\mathbb F_{64}\)，但不包含 \(\mathbb F_{16}\)；\(\mathbb F_4\cap\mathbb F_8=\mathbb F_2\)，而 \(\mathbb F_4\mathbb F_8=\mathbb F_{64}\)。不能仅凭域的元素个数较小就断定存在包含关系。

\ProseParagraphBreak
这也让我们能直接检验域生成与乘法群生成的区别。取 \(\mathbb F_{16}^\times\) 的一个生成元 \(g\)，其阶为十五，故 \(g^3\) 的阶为五。\(\mathbb F_{16}\) 的真子域只有 \(\mathbb F_2\) 和 \(\mathbb F_4\)，它们的乘法群阶分别为一和三，都不可能含有五阶元。所以 \(g^3\) 不在任何真子域中，从而 \(\mathbb F_2(g^3)=\mathbb F_{16}\)；但 \(g^3\) 只生成五阶乘法子群，仍不是有限域的乘法本原元。

\subsection{乘法群生成元的计算}
\label{b1:subsec:1702:04}

设 \(N=q-1\)，且已分解出 \(N\) 的不同素因子 \(\ell_1,\ldots,\ell_s\)。元素 \(g\in\mathbb F_q^\times\) 是本原元当且仅当
\begin{equation}\label{b1:eq:primitive-finite-field-test}
g^{N/\ell_j}\ne1\qquad(1\leq j\leq s).
\end{equation}
事实上，其阶 \(d\) 整除 \(N\)；若 \(d<N\)，则 \(N/d>1\)，取 \(N/d\) 的任意一个素因子 \(\ell_j\)，便有 \(d\mid N/\ell_j\)，所以 \(g^{N/\ell_j}=1\)，违反检验。这个 \(\ell_j\) 必也是 \(N\) 的素因子，故检查列出的不同素因子已能发现每一种 \(d<N\) 的情形，不必再检查 \(\ell_j^2,\ell_j^3\) 等素数幂。反过来，若 \(d=N\)，因为 \(N/\ell_j<N\)，这些幂都不可能为一。若 \(N=1\)，没有素因子，唯一的非零元就是一，空的检验也与结论相符。

\ProseParagraphBreak
在 \(\mathbb F_9=\mathbb F_3(u)\)、\(u^2=2\) 中，取 \(g=1+u\)。直接计算
\[
g^2=2u,\qquad g^4=2,\qquad g^8=1.
\]
八的唯一素因子为二，而 \(g^4\ne1\)，所以 \(g\) 是本原元。这也给出一个实际搜索办法：逐个枚举非零元素，用重复平方计算上述幂，直到找到全部检验通过的元素。定理保证搜索终会成功；若不知道 \(q-1\) 的素因子，还须先处理整数分解，不能把这一步的代价忽略。
